\answer
\begin{enumerate}

%% 10章 問1（10.2節）
\item[10.2節 問1]
全波整流後の平滑キャパシタは波形のピーク値まで充電されるから，
無負荷のDCリンク電圧はほぼ
$\sqrt{2} \times 200 \approx 283$\,Vになる。
負荷をつなぐと充放電のリプルの分だけこれより下がる（8章）。

%% 10章 問2（10.3節）
\item[10.3節 問2]
$\alpha$の単位ラジアンは「弧の長さ$\div$半径」で定義される無次元量である。
したがって$\alpha/\pi$も$\sin 2\alpha/(2\pi)$も無次元で，根号の中全体が無次元になる。
残るのは$V_0/\sqrt{2}$で，$V_0$は電圧だから右辺の単位はVである。左辺（電圧の実効値）と一致する。

%% 10章 問3（10.3節）
\item[10.3節 問3]
式\ref{eq10.3}の根号の中を計算する。
$\alpha = \pi/4$：$1 - 1/4 + \sin(\pi/2)/(2\pi) = 0.75 + 0.159 \approx 0.909$。
実効値は$\sqrt{0.909} \approx 0.95$倍，電力は$0.91$倍。
$\alpha = 3\pi/4$：$1 - 3/4 + \sin(3\pi/2)/(2\pi) = 0.25 - 0.159 \approx 0.091$。
実効値は$\sqrt{0.091} \approx 0.30$倍，電力は$0.091$倍。
点弧角が小さいうちは電力があまり減らず，大きくなると急に絞られることが\ref{fig10.3}からも読み取れる。

%% 10章 問4（10.4節）
\item[10.4節 問4]
サイクロコンバータ：三相ブリッジ1つにサイリスタ6個。
1つの出力相につき正群・負群の2ブリッジで$6 \times 2 = 12$個，
3相出力では$12 \times 3 = 36$個のサイリスタが必要である。
マトリックスコンバータ：双方向スイッチは$3 \times 3 = 9$個で，
IGBTは$9 \times 2 = 18$個（ダイオードも18個）で済む。
素子数の少なさもマトリックスコンバータの魅力の1つである。
\end{enumerate}
